\documentclass{article}
\usepackage{amsmath}
\usepackage{amsthm}
\usepackage{amssymb}
\usepackage{setspace}
\theoremstyle{definition}
\newtheorem{definition}{Def}
\newtheorem{theorem}[definition]{Thm}
\newtheorem{corollary}[definition]{Cor}
\newtheorem{proposition}[definition]{Prop}
\newtheorem{remarks}[definition]{Rmk}
\newtheorem{notation}[definition]{Notation}
\onehalfspacing

\title{Chapter1}
\author{Ricky Dai}
\date{September 2026}

\begin{document}

\maketitle

\section{Introduction}
\begin{definition}
    If $A$ is any set \, (whose elements may be numbers or any other objects), we write $x \in A$ to indicate that $x$ is a member \, (or an element) of $A$. If $x$ is not a member of $A$, we write $x \notin A$.\par
    The set which contains no element will be called the \textit{empty set}. If a set has at least one element, it is called \textit{nonempty}\par
    If $A$ and $B$ are sets, and if every element of $A$ is an element of $B$, we say that $A$ is a subset of $B$, and we write $A \subset B$, or $B \supset A$. If, in addition, there is an element of $B$ which is not in $A$, then $A$ is said to be a \textit{proper} subset of $B$. Note that $A \subset A$ for every set A.\par
    If $A \subset B$ and $B \subset A$, we write $A = B$. Otherwise, $A \neq B$.
\end{definition}

\begin{definition}
    The set of all rational numbers will be denoted by $\mathbb{Q}$.
\end{definition}

\section{Ordered Sets}
\begin{definition}
    Let $S$ be a set. An \textit{order} on S is a relation, denoted by $<$, with the following two properties \par
    (i) If $x \in S$ and $y \in S$, then one and only one of the statements
    \[
    x < y,\;x = y,\;y < x
    \]\par
    is true. \par
    (ii) If $x, y, z \in S$, if $x < y$ and $y < z$, then $x < z$.
\end{definition}

\begin{definition}
    An \textit{ordered set} is a set $S$ in which an order is defined.
\end{definition}
\begin{definition}
    Suppose $S$ is an ordered set, and $E \subset S$. If there exists a $\beta \in S$ such that $x \leq \beta$ for every $x \in E$, we say that E is \textit{bounded above}, and call $\beta$ an \textit{upper bound} of $E$.
\end{definition}
\begin{definition}
    Suppose $S$ is an ordered set, $E \subset S$, and $E$ is bounded above. Suppose there exists an $\alpha \in E$ with the following properties:\par
    (i) $\alpha$ is an upper bound of $E$.\par
    (ii) If $\gamma < \alpha$ then $\gamma$ is not an upper bound of $E$.\\
    Then $\alpha$ is called the \textit{least upper bound} of $E$ [that there is at most one such $\alpha$ is clear from (ii)] or the \textit{supremum} of $E$, and we write
    \[
    \alpha = sup\;E.
    \]\par
    The \textit{greatest lower bound}, or \textit{infimum}, of a set $E$ which is bounded below is defined in the same manner: The statement
    \[
    \alpha = inf\;E
    \]
    means that $\alpha$ is a lower bound of $E$ and that no $\beta$ with $\beta > \alpha$ is a lower bound of $E$.
\end{definition}
\begin{definition}
    An ordered set S is said to have the \textit{least-upper-bound property} if the following is true:\par
    If $E \subset S$, $E$ is not empty, and $E$ is bounded above, then sup $E$ exists in $S$.
\end{definition}
\begin{theorem}
    Suppose $S$ is an ordered set with the least-upper-bound property, $B \subset S$, $B$ is not empty, and $B$ is bounded below. Let $L$ be the set of all lower bounds of $B$. Then
    \[
    \alpha = sup\;L
    \]
    exists in $S$ and $\alpha = inf\;B$.\par
    In particular, $inf\;B$ exists in $S$.
\end{theorem}

\section{Fields}
\begin{definition}
    A \textit{field} is a set $\mathbb{F}$ with two operations, called \textit{addition} and \textit{multiplication}, which satisfy the following so-called "field axioms" (A), (M), and (D):\par
    \textbf{(A) Axioms for addition}\par
    \hspace*{2em}(A1) $x \in \mathbb{F}$ and $y \in \mathbb{F}$, then their sum $x+y$ is in $\mathbb{F}$.\par
    \hspace*{2em}(A2) Addition is commutative: $x+y=y+x$ for all $x, y \in \mathbb{F}$. \par
    \hspace*{2em}(A3) Addition is associative: $(x+y)+z=x+(y+z)$ for all $x, y, z \in \mathbb{F}$.\par
    \hspace*{2em}(A4) $\mathbb{F}$ contains an element $0$ such that $0+x=x$ for every $x \in \mathbb{F}$.\par
    \hspace*{2em}(A5) To every $x \in \mathbb{F}$ corresponds an element $-x \in \mathbb{F}$ such that 
    \[
    x+(-x)=0.
    \]\par
    \textbf{(M) Axioms for multiplication}\par
    \hspace*{2em}(M1) If $x,y \in \mathbb{F}$, then their product $xy$ is in $\mathbb{F}$.\par
    \hspace*{2em}(M2) Multiplication is commutative: $xy=yx$ for all $x, y \in \mathbb{F}$.\par
    \hspace*{2em}(M3) Multiplication is associative: $(xy)z=x(yz)$ for all $x, y, z \in \mathbb{F}$.\par
    \hspace*{2em}(M4) $\mathbb{F}$ contains an element $1\neq0$ such that $1x=x$ for every $x \in \mathbb{F}$.\par
    \hspace*{2em}(M5) If $x \in \mathbb{F}$ and $x \neq 0$ then there exists an element $1/x \in \mathbb{F}$ such that 
    \[
    x \cdot (1/x)=1.
    \]\par
    \textbf{(D) The distribution law}
    \[
    x(y+z)=xy+xz
    \]\par
    \hspace*{2em}holds for all $x, y, z\in \mathbb{F}$.
\end{definition}
\begin{proposition}
    The axioms for addition imply the following statements.\par
    (a) If $x+y=x+z$ then $y=z$.\par
    (b) If $x+y=x$ then $y=0$.\par
    (c) If $x+y=0$ then $y=-x$.\par
    (d) $-(-x)=x.$
\end{proposition}
\begin{proposition}
    The axioms for multiplication imply the following statements\par
    (a) If $x\neq0$ and $xy=yz$ then $y=z$.\par
    (b) If $x\neq0$ and $xy=x$ then $y=1$.\par
    (c) If $x\neq0$ and $xy=1$ then $y=1/x$.\par
    (d) If $x\neq0$ then $1/(1/x))=x$.
\end{proposition}
\begin{proposition}
    The field axioms imply the following statements, for any $x, y, z \in \mathbb{F}$.\par
    (a) $0x=0$.\par
    (b) If $x\neq0$ and $y\neq0$ then $xy\neq0$.\par
    (c) $(-x)y=-(xy)=x(-y)$.\par
    (d) $(-x)(-y)=xy$.\par
\end{proposition}
\begin{definition}
    An \textit{ordered field} is a \textit{field} $\mathbb{F}$ which is also an \textit{ordered set}, such that\par
    (i) $x+y<x+z$ if $x, y, z \in \mathbb{F}$ and $y<z$,\par
    (ii) $xy>0$ if $x \in \mathbb{F}, y\in\mathbb{F}$, and $y>0$.\\
    If $x>0$, we call $x$ \textit{positive}; if $x<0$, $x$ is \textit{negative}.
\end{definition}
\begin{proposition}
    The following statements are true in every ordered field.\par
    (a) If $x>0$ then $-x<0$, and vice versa.\par
    (b) If $x>0$ and $y < z$ then $xy < xz$.\par
    (c) If $x<0$ and $y < z$ then $xy>xz$.\par
    (d) If $x\neq0$ then $x^2>0$. In particular, $1>0$.\par
    (e) If $0<x<y$ then $0<1/y<1/x$.
\end{proposition}

\section{The Real Field}
\begin{theorem}
    There exists an ordered field $\mathbb{R}$ which has the least-upper-bound property. Moreover, $\mathbb{R}$ contains $\mathbb{Q}$ as a subfield.
\end{theorem}
\begin{theorem}
    \leavevmode\par
    (a) If $x\in \mathbb{R}$, $y\in \mathbb{R}$, and $x>0$, then there is a positive integer $n$ such that
    \[
    nx > y.
    \]\par
    (b) If $x\in \mathbb{R}$, $y\in\mathbb{R}$, and $x<y$, then there exists a $p \in Q$ such that $x<p<y$.
\end{theorem}
\begin{theorem}
    For every real $x>0$ and every integer $n>0$, there is one and only one positive real $y$ such that $y^n=x$.\par
    This number $y$ is written $\sqrt[n]{x}$ or $x^{1/n}$.
\end{theorem}
\begin{corollary}
    If $a$ and $b$ are positive real numbers and $n$ is a positive integer, then
    \[
    (ab)^{1/n}=a^{1/n}b^{1/n}.
    \]
\end{corollary}

\section{The Extended Real Number System}
\begin{definition}
    The extended real number system consists of real field $\mathbb{R}$ and two symbols, $+\infty$ and $-\infty$. We preserve the original order in $\mathbb{R}$ and define
    \[
    -\infty<x<+\infty
    \]
    for every $x\in\mathbb{R}$.
\end{definition}

\section{The Complex Field}
\begin{definition}
    A complex number is an ordered pair $(a, b)$ of real numbers. "Ordered" means that $(a, b)$ and $(b, a)$ are regarded as distinct if $a\neq b$.\par
    Let $x=(a, b)$, $y=(c, d)$ be two complex numbers, we write $x=y$ if and only if $a=c$ and $b=d$.We define:
    \[
    \begin{aligned}
        x+y&=(a+c, b+d) \\
        xy&=(ac-bd, ad+bc)
    \end{aligned}
    \]
\end{definition}
\begin{theorem}
    These definitions of addition and multiplication turn the set of all complex numbers into a field, which $(0, 0)$ and $(1, 0)$ in the role of $0$ and $1$.
\end{theorem}
\begin{theorem}
    For any real numbers $a$ and $b$ we have
    \[
    (a, 0) + (b, 0) = (a+b, 0),\qquad (a, 0)(b, 0)= (ab, 0)
    \]
\end{theorem}
\begin{definition}
    $i=(0, 1)$.
\end{definition}
\begin{theorem}
    $i^2=-1$.
\end{theorem}
\begin{theorem}
    If $a$ and $b$ are real, then $(a, b)=a+bi$.
\end{theorem}
\begin{definition}
    If $a$, $b$ are real and $z = a+bi$, then the complex number $\overline{z}=a-bi$ is called the \textit{conjugate} of $z$. The numbers $a$ and $b$ are the \textit{real part} and the \textit{imaginary part} of $z$, respectively.\par
    We shall occasionally write
    \[
    a=Re(z),\qquad b=Im(z).
    \]
\end{definition}
\begin{theorem}
    If $z$ and $w$ are complex, then\par
    (a) $\overline{z+w}=\overline{z}+\overline{w}$,\par
    (b) $\overline{zw}=\overline{z}\cdot\overline{w}$,\par
    (c) $z+\overline{z}=2Re(z),z-\overline{z}=2iIm(z)$,\par
    (d) $z\overline{z}$ is real and positive(except when $z=0$).
\end{theorem}
\begin{definition}
    If $z$ is a complex number, its \textit{absolute value} $|z|$ is the non-negative square root of $z\overline{z}$; that is $|z|=(z\overline{z})^{1/2}$.
\end{definition}
\begin{theorem}
    Let $z$ and $w$ be complex numbers. Then\par
    (a) $|z|>0$ unless $z=0$, $|0|=0$,\par
    (b) $|\overline{z}|=|z|$,\par
    (c) $|zw|=|z||w|$,\par
    (d) $|Re(z)|\leq|z|$,\par
    (e) $|z+w|\leq|z|+|w|$.
\end{theorem}
\begin{notation}
    If $x_1,...,x_n$ are complex numbers, we write
    \[
    x_1+x_2+\cdot\cdot\cdot+x_n=\sum_{j=1}^n x_j.
    \]
\end{notation}
\begin{theorem}
    If $a_1,...,a_n$ and $b_1,...,b_n$ are complex numbers, then
    \[
    |\sum_{j=1}^n a_j\overline{b_j}|^2 \leq \sum_{j=1}^n |a_j|^2 \sum_{j=1}^n |b_j|^2.
    \]
\end{theorem}

\section{Euclidean Spaces}
\begin{definition}
    For each positive integer $k$, let $\mathbb{R}^k$ be the set of all ordered k-tuples 
    \[
    \textbf{x} = (x_1, x_2, ..., x_k),
    \]
    where $x_1,...,x_k$ are real numbers, called the \textit{coordinates} of $\textbf{x}$. The elements of $\mathbb{R}^k$ are called points, or vectors, especially when $k > 1$. We shall denote vectors by boldface letters. If $\textbf{y}=(y_1,...,y_k)$ and if $\alpha$ is a real number, put
    \[
    \begin{aligned}
        \textbf{x}+\textbf{y} &= (x_1+y_1,...,x_k+y_k)\\
        \alpha \textbf{x} &= (\alpha x_1,...,\alpha x_n)
    \end{aligned}
    \]
    so that $\textbf{x}+\textbf{y}\in\mathbb{R}^k$ and $\alpha\textbf{x}\in\mathbb{R}^k$. This defines addition of vectors, as well as multiplication of a vector by a real number(a scalar). These two operations satisfy the commutative, associative, and distribution laws and make $\mathbb{R}^k$ into a \textit{vector space over the real field}. The zero element of $\mathbb{R}^k$ (sometimes called the \textit{origin} or the \textit{null vector}) is the point $\textbf{0}$, all of whose coordinates are $0$.\par
    We also define the so-called "inner-product" (or scalar product) of $\textbf{x}$ and $\textbf{y}$ by
    \[
    \textbf{x}\cdot\textbf{y}=\sum_{i=1}^k x_i y_i
    \]
    and the \textit{norm} of $\textbf{x}$ by
    \[
    |\textbf{x}| = (\textbf{x}\cdot\textbf{x})^{1/2}=(\sum_{i=1}^k x_i^2)^{1/2}.
    \]
    The structure now defined (the vector space $\mathbb{R}^k$ with the above inner product and norm) is called euclidean k-space.
\end{definition}
\begin{theorem}
    Suppose $\textbf{x},\textbf{y},\textbf{z}\in\mathbb{R}^k$, and $\alpha$ is real. Then \par
    (a) $|\textbf{x}| \ge 0$;\par
    (b) $|\textbf{x}|=0$ if and only if $\textbf{x}=0$;\par
    (c) $|\alpha\textbf{x}|=|\alpha||\textbf{x}|$;\par
    (d) $|\textbf{x}\cdot\textbf{y}|\leq|\textbf{x}||\textbf{y}|$;\par
    (e) $|\textbf{x}+\textbf{y}|\leq|\textbf{x}|+|\textbf{y}|$;\par
    (f) $|\textbf{x}-\textbf{z}|\leq|\textbf{x}-\textbf{y}|+|\textbf{y}-\textbf{z}|$.
\end{theorem}
\begin{remarks}
    (a), (b), and (f) will allow us to regard $\mathbb{R}^k$ as a metric space.
\end{remarks}
\end{document}
